\chapter{Introduction}\label{ch:intro}

\section{The basic definition}

Let $\cC$ be a category possessing an initial object. The central object of
study in this monograph is the following.

\begin{definition*}[informal version of \cref{def:cat}]
A \emph{$\cC$-cat} is a tuple $K=(H,T,(L_\alpha)_{\alpha\in A})$ in which $H$
is an initial object of $\cC$, $T$ is a terminal object of $\cC$, $A$ is a
small set, and each $L_\alpha$ is an object of $\cC$, together with the
canonical morphisms
\[
\begin{tikzcd}
  & L_\alpha \arrow[dr, "t_\alpha"] & \\
  H \arrow[ur, "h_\alpha"] \arrow[rr, "\sigma"'] & & T
\end{tikzcd}
\]
for $\alpha\in A$.
\end{definition*}

The object $H$ is the \emph{head}, $T$ the \emph{tail}, and the $L_\alpha$
the \emph{legs} of $K$; the morphism $\sigma$ is the \emph{spine}. A morphism
of $\cC$-cats consists of a reindexing function between the sets of legs and
morphisms between corresponding heads, legs and tails, subject to the
commutativity of a family of squares (\cref{def:morphism}). The resulting
category is denoted $\Fel(\cC)$.

\section{An immediate reduction}

The first substantive observation of the theory, which we record here so that
it is not obscured by what follows, is that the head and tail of a cat carry no
information.

\begin{maintheorem}[Fundamental Theorem of Cat Theory; \cref{thm:fundamental}]
Let $\cC$ be a category with an initial object. Every commutativity condition
in the definition of a morphism of $\cC$-cats is automatically satisfied. The
functor $\Fel(\cC)\to\Fam(\cC)$ which forgets the head and the tail is fully
faithful, and it is an equivalence of categories if and only if $\cC$ has a
terminal object. It is an isomorphism of categories if and only if $\cC$ has
exactly one initial object and exactly one terminal object.
\end{maintheorem}

Here $\Fam(\cC)$ is the category of small families of objects of $\cC$, which
is, by a classical observation, the free completion of $\cC$ under small
coproducts. The Fundamental Theorem converts every subsequent question about
$\cC$-cats into a question about families. We have nevertheless chosen to
state and prove all results in feline terms, partly in order to verify that
the feline formalism is internally consistent and partly because certain
features of families (for instance, the r\^ole played by the initial object
of $\cC$ in the existence of limits) become visible only when the head is
recorded explicitly.

\section{Principal results}

We list the principal results. The labels refer to the body of the text.

\begin{mainthm}[Limits and colimits; \cref{thm:limits,thm:colimits}]\label{mthm:limits}
Let $\cC$ have an initial and a terminal object. For every small category
$\cJ$, the category $\Fel(\cC)$ has all limits of shape $\cJ$ if and only if
$\cC$ does. The category $\Fel(\cC)$ has all small colimits if and only if
$\cC$ has all small connected colimits. The hypothesis that $\cC$ have an
initial object is essential for the first statement.
\end{mainthm}

\begin{mainthm}[The feline adjoint string; \cref{thm:adjoint-string,prop:Delta0-left,prop:Gamma-right}]\label{mthm:string}
The leg functor $\legs\colon\Fel(\cC)\to\Set$ sits in a string of adjunctions
\[
\Delta_0\dashv \legs \dashv \Delta_1 \dashv \points,
\]
where $\Delta_0$ (resp.\ $\Delta_1$) sends a set $S$ to the cat with $S$
legs, each of which is a head (resp.\ a tail), and $\points$ sends a cat to
the set of global elements of its legs. The functor $\Delta_0$ has a left
adjoint if and only if $\cC$ is pointed, and $\points$ has a right adjoint
if and only if $\cC$ is pointed or the global-sections functor
$\cC(1,-)\colon\cC\to\Set$ has a right adjoint. For $\cC=\Set$ the string
extends to five functors and cannot be extended further.
\end{mainthm}

\begin{mainthm}[Feline Yoneda; \cref{thm:yoneda-ff,thm:yoneda-image,thm:head-tail-eval}]\label{mthm:yoneda}
The functor $\Yon\colon\Fel(\cC)\to\PSh(\cC)$ sending a cat to the coproduct
of the presheaves represented by its legs is fully faithful. Its essential
image consists of those presheaves whose category of elements is a disjoint
union of categories with terminal objects. Under $\Yon$, the leg functor is
evaluation at the head and the points functor is evaluation at the tail.
\end{mainthm}

\begin{mainthm}[Reconstruction; \cref{thm:reconstruction}]\label{mthm:reconstruction}
If $\Fel(\cC)\eqv\Fel(\cD)$, then $\cC\eqv\cD$.
\end{mainthm}

\begin{mainthm}[Computations; \cref{thm:Fel-Set-J,cor:Fel-presheaf,thm:iterated-point}]\label{mthm:computations}
For every small category $\cJ$ there is an equivalence
$\Fel(\Set^{\cJ})\eqv\Set^{\cJ^{\triangleright}}$, where $\cJ^{\triangleright}$
is $\cJ$ with a terminal object adjoined. Consequently
$\Fel(\PSh(\cD))\eqv\PSh(\cD^{\triangleleft})$ for every small $\cD$, the
$n$-fold iterated category of cats over the terminal category is equivalent to
$\Set^{[n-1]}$, and in all cases the leg functor corresponds to evaluation at
the adjoined object.
\end{mainthm}

\begin{mainthm}[Feline toposes; \cref{thm:topos}]\label{mthm:topos}
If $\cE$ is an elementary topos with small products, then $\Fel(\cE)$ is an
elementary topos. Its subobject classifier is a two-legged cat whose legs are
the subobject classifier of $\cE$ and the terminal object.
\end{mainthm}

\begin{mainthm}[CMG Correspondence; \cref{thm:cmg}]\label{mthm:cmg}
Let $X$ be an object of $\cC$. Every Gideon Decomposition of $X$ satisfies the
Chloe Condition and is Margaret compatible. Conversely, if $X$ admits at least
one Gideon Decomposition, then a cocone from a cat satisfying the Chloe
Condition to $X$ is a Gideon Decomposition if and only if it is Margaret
compatible; Gideon Decompositions of $X$ are unique up to unique isomorphism.
Neither hypothesis can be omitted.
\end{mainthm}

\begin{mainthm}[Recognition; \cref{thm:recognition}]\label{mthm:recognition}
A locally small category $\cE$ is equivalent to $\Fel(\cD)$ for some category
$\cD$ with an initial and a terminal object if and only if $\cE$ has small
coproducts, every object of $\cE$ admits a Gideon Decomposition, and the full
subcategory of connected objects of $\cE$ (equivalently, of objects satisfying
the Chloe Condition) has an initial object. The existence of a terminal object
in $\cE$ is implied by these conditions, a Gideon Decomposition being a cat and
therefore having a tail.
\end{mainthm}

\section{The CMG framework}

Three notions recur in the later chapters. The \emph{Chloe Condition}
(\cref{sec:chloe}) is a condition on a cat: that each of its legs be a connected object of $\cC$.
\emph{Margaret Compatibility} is a condition on a cocone from the legs of a
cat to an object $X$: that every morphism from a connected object into $X$
factor uniquely through exactly one leg (\cref{sec:margaret}). A
\emph{Gideon Decomposition} of $X$ (\cref{sec:gideon}) is a cat satisfying the Chloe Condition together with a coproduct cocone
from its legs to $X$. The CMG Correspondence Theorem (\cref{mthm:cmg}),
proved in \cref{ch:cmg}, shows
that, in the presence of at least one Gideon Decomposition, Margaret
Compatibility characterizes Gideon Decompositions among cocones from cats
satisfying the Chloe Condition, and that Gideon Decompositions are essentially
unique. We identify these notions with familiar ones in
\cref{rem:cmg-honest}: the Chloe Condition is connectedness of legs, Margaret
Compatibility is the invertibility of a comparison of restricted Yoneda
presheaves, and Gideon Decompositions are decompositions into connected
components. The three notions are introduced where they arise (in
\cref{ch:anatomy,ch:yoneda,ch:representation} respectively), recur in
\cref{ch:limits,ch:adjunctions,ch:monoidal,ch:fibrations}, and are brought
together in \cref{ch:cmg}.

\section{Relation to the literature}

We make no claim of novelty. The category of small families $\Fam(\cC)$, its
description as the free coproduct completion, the family fibration
$\Fam(\cC)\to\Set$, and the elementary facts concerning its limits,
colimits, connected objects and Yoneda embedding are classical; general
references are \cite{Borceux1994,Jacobs1999,Johnstone2002,MacLane1998,Riehl2016}.
The characterization of coproduct completions by connected objects, which
underlies \cref{mthm:recognition}, is likewise well known and elementary. We have included
complete proofs because the feline formulation of each statement differs
superficially from the classical one and because we wished the account to be
self-contained. Where we have not located a statement in the literature in
precisely the form given, we have nevertheless refrained from describing it as
new: in every such case the statement is an elementary consequence of the
Fundamental Theorem together with well-known facts.

The terminology of this monograph (heads, tails, legs, kittens, the Chloe
Condition, Margaret Compatibility, Gideon Decompositions, and so on) is new,
but new terminology does not constitute new mathematics. The reader is
referred to \cref{app:dictionary} for a translation of every feline term into
standard language.

\section{Conventions}

Our set-theoretic and categorical conventions are fixed in \cref{sec:conventions}.
We write composition of morphisms in applicative order, $g\circ f$ or $gf$
for ``first $f$, then $g$''. Diagrams are said to commute if all composites
of parallel paths between any two vertices agree. Throughout, \emph{the}
category $\cC$ is a fixed ambient category; the standing hypotheses on it are
stated at the beginning of each chapter.
